\documentclass[12pt,a4paper]{article}
\usepackage{fontspec}
\setmainfont{Times New Roman}
\usepackage[top=2.0cm,bottom=2.0cm,left=2.5cm,right=2.0cm]{geometry}
\usepackage{enumitem}
\usepackage{tabularx}
\usepackage{xltabular}
\usepackage{booktabs}
\usepackage{array}
\newcolumntype{Y}{>{\centering\arraybackslash}X}
\newcolumntype{C}[1]{>{\centering\arraybackslash}p{#1}}
\newcolumntype{L}[1]{>{\raggedright\arraybackslash}p{#1}}
\newcolumntype{R}[1]{>{\raggedleft\arraybackslash}p{#1}}
\pagestyle{plain}
\linespread{1.15}
\setlength{\parskip}{4pt plus 1pt minus 1pt}
\sloppy

\begin{document}
% --- HEADER 2 CỘT CHUẨN NGHỊ ĐỊNH 30/2020/NĐ-CP ---
\noindent
\begin{minipage}[t]{0.46\textwidth}
    \begin{center}
        \fontsize{11pt}{13pt}\selectfont
        CÔNG AN THÀNH PHỐ HÀ NỘI\\
        \textbf{PHÒNG CẢNH SÁT HÌNH SỰ (PC02)}\\[-0.2em]
        \rule{3.2cm}{0.6pt}\\[0.25cm]
        \fontsize{11pt}{13pt}\selectfont Số: \textbf{04/BC-XMNT}
    \end{center}
\end{minipage}
\hfill
\begin{minipage}[t]{0.52\textwidth}
    \begin{center}
        \fontsize{11pt}{13pt}\selectfont
        \textbf{CỘNG HÒA XÃ HỘI CHỦ NGHĨA VIỆT NAM}\\
        \textbf{Độc lập - Tự do - Hạnh phúc}\\[-0.2em]
        \rule{3.6cm}{0.6pt}\\[0.25cm]
        \textit{Hà Nội, ngày 25 tháng 07 năm 2016}
    \end{center}
\end{minipage}

\vspace{0.5cm}

\begin{center}
    \textbf{\Large BÁO CÁO XÁC MINH NHÂN THÂN, LAI LỊCH CỦA NẠN NHÂN}\\[0.15cm]
    \textit{(Nạn nhân: Nguyễn Văn Khang — Vụ án mạng số 14 Đường Bờ Sông)}
\end{center}

\vspace{0.3cm}



\vspace{0.3cm}\noindent\textbf{\large I. LÝ LỊCH TƯ PHÁP \& QUAN HỆ GIA ĐÌNH}\\[0.15cm]

\begin{itemize}[leftmargin=*, itemsep=2pt, topsep=2pt, parsep=0pt]
  \item \textbf{Họ và tên:} \textbf{NGUYỄN VĂN KHANG}
  \item \textbf{Sinh ngày:} 10/11/1988 tại Hà Nội (28 tuổi).
  \item \textbf{Số CCCD/CMND:} \texttt{830514692}.
  \item \textbf{Trình độ học vấn:} 12/12.
  \item \textbf{Tình trạng gia đình:}
\begin{itemize}[leftmargin=*, itemsep=2pt, topsep=2pt, parsep=0pt]
  \item Bố mẹ đẻ mất sớm do tai nạn giao thông năm 2010.
  \item Sống cùng ông nội (ông Nguyễn Văn Thọ) tại số 14 Đường Bờ Sông. Ông Thọ vừa qua đời cách đây 3 tháng.
  \item Họ hàng gần nhất: Gia đình chú ruột trú tại Phường Phân khu Cảng.


\end{itemize}
\end{itemize}

\vspace{0.3cm}\noindent\textbf{\large II. HOẠT ĐỘNG KINH TẾ \& QUAN HỆ XÃ HỘI}\\[0.15cm]

\begin{enumerate}[label=\arabic*., leftmargin=*, itemsep=2pt, topsep=2pt, parsep=0pt]
  \item \textbf{Nghề nghiệp chính:} Không có công ăn việc làm ổn định. Mở dịch vụ "Tài chính tiêu dùng cá nhân" (thực chất là cho vay nặng lãi).
  \item \textbf{Tiền án, tiền sự:}
\begin{itemize}[leftmargin=*, itemsep=2pt, topsep=2pt, parsep=0pt]
  \item 01 tiền sự năm 2012 về hành vi \textit{"Gây rối trật tự công cộng"} (đánh nhau tại quán bia).
\end{itemize}
  \item \textbf{Đánh giá nguy cơ mâu thuẫn:}
\begin{itemize}[leftmargin=*, itemsep=2pt, topsep=2pt, parsep=0pt]
  \item Do đặc thù nghề nghiệp cho vay nặng lãi, nạn nhân có quan hệ xã hội phức tạp, nhiều đối tượng có hiềm khích nợ nần hoặc tranh chấp tài sản.
\end{itemize}
\end{enumerate}
\vspace{0.8cm}
\noindent
\hfill
\begin{minipage}[t]{0.50\textwidth}
    \begin{center}
        \textbf{CÁN BỘ XÁC MINH}\\[0.1cm]
        \textit{(Ký, ghi rõ họ tên)}\\[1.6cm]
        \textbf{Trung úy Lê Hoàng Long}
    \end{center}
\end{minipage}
\end{document}
